% GENERATED FILE. Edit canonical lessons, then run npm run generate:books.
% canonical-source-hash: fnv1a64:945fc75b2c21e5ee
% canonical-lessons: LA-C477-pecus, LA-C477-grex, LA-C477-armentum, LA-C477-pasco, LA-C477-pascuum

\chapter{Livestock, Flock, Herd, To pasture, Pasture}
\label{ch:la-livestock-flock-herd-to-pasture-pasture}
\hlchaptermodality{477}

\begin{chapteropening}
\textbf{By the end of this chapter:} \emph{I can talk about farm animals and their grazing: livestock, a flock, a herd, putting animals out to pasture, and the pasture.}

Complete the last lesson of chapter 477: I can talk about farm animals and their grazing: livestock, a flock, a herd, putting animals out to pasture, and the pasture.
\end{chapteropening}

\section[pecus, pecoris]{pecus, pecoris — livestock, cattle, a herd of animals}

\label{lesson:LA-C477-pecus}



\begin{quote}
Before the new one: say the Latin for a stalk, then the Latin for a granary.
\end{quote}

\subsection*{You'll want to know: pecus, pecoris}
\textbf{pecus, pecoris} — \textquotedblleft{}livestock, cattle, a herd of animals\textquotedblright{}.

\begin{quote}
Agricola multum pecus habet: ovēs, caprās, bovēs.
\end{quote}

\begin{quote}
\emph{The farmer has a lot of livestock: sheep, goats, oxen.}
\end{quote}

Third declension, neuter: \emph{pecus, pecoris n.}, the animals as a group. Another word \emph{pecus, pecudis f.} is one beast. \emph{Pecūnia} 'money' comes from the same root: wealth was once counted in cattle.

\begin{cousinweb}
\emph{pecus} is inherited from PIE \emph{pe\'{k}u-} 'livestock' (EtymDB); Lewis \& Short list Gothic \emph{faihu}, German \emph{Vieh} and English \emph{fee} as cousins. \emph{pecūnia} 'money' is built on the same root (Lewis \& Short).
\end{cousinweb}

\begin{grammarlens}[title={What you've built}]
One more word for telling, explaining and describing at B1.
\end{grammarlens}

\subsection*{Your turn}
\begin{itemize}
\raggedright
  \item \emph{Say it:} \emph{pecus, pecoris}
  \item \emph{Say it:} \emph{pecus, pecoris}, once more
  \item \emph{Say it:} a sentence of your own with \emph{pecus, pecoris}
  \item \emph{Recall:} say the Latin for a stalk, then the Latin for a granary, then say \emph{pecus, pecoris} again
\end{itemize}

\begin{tcolorbox}[breakable,colback=teal!4,colframe=teal!35!black,title={Before you move on}]
What does pecus, pecoris mean? (\textquotedblleft{}Livestock, cattle, a herd of animals\textquotedblright{}.) Say it once more.
\end{tcolorbox}
\section[grex, gregis]{grex, gregis — a flock, a herd; a band of people}

\label{lesson:LA-C477-grex}



\begin{quote}
Before the new one: say the Latin for a granary, then the Latin for livestock.
\end{quote}

\subsection*{You'll want to know: grex, gregis}
\textbf{grex, gregis} — \textquotedblleft{}a flock, a herd; a band of people\textquotedblright{}.

\begin{quote}
Pastor gregem ovium ad flūmen dūcit.
\end{quote}

\begin{quote}
\emph{The shepherd leads the flock of sheep to the river.}
\end{quote}

Third declension, masculine: \emph{grex, gregis m.} Also a group of people. English \emph{gregarious}, \emph{congregate} and \emph{egregious} (out of the flock) come from it.

\begin{cousinweb}
\emph{grex} is inherited from PIE \emph{h\textsubscript{2}ger-} 'to assemble, to gather' (EtymDB). English \emph{gregarious}, \emph{congregate} and \emph{egregious} come from it.
\end{cousinweb}

\begin{grammarlens}[title={What you've built}]
One more word for telling, explaining and describing at B1.
\end{grammarlens}

\subsection*{Your turn}
\begin{itemize}
\raggedright
  \item \emph{Say it:} \emph{grex, gregis}
  \item \emph{Say it:} \emph{grex, gregis}, once more
  \item \emph{Say it:} a sentence of your own with \emph{grex, gregis}
  \item \emph{Recall:} say the Latin for a granary, then the Latin for livestock, then say \emph{grex, gregis} again
\end{itemize}

\begin{tcolorbox}[breakable,colback=teal!4,colframe=teal!35!black,title={Before you move on}]
What does grex, gregis mean? (\textquotedblleft{}A flock, a herd; a band of people\textquotedblright{}.) Say it once more.
\end{tcolorbox}
\section[armentum, armentī]{armentum, armentī — a herd (of cattle or horses)}

\label{lesson:LA-C477-armentum}



\begin{quote}
Before the new one: say the Latin for livestock, then the Latin for a flock.
\end{quote}

\subsection*{You'll want to know: armentum, armentī}
\textbf{armentum, armentī} — \textquotedblleft{}a herd (of cattle or horses)\textquotedblright{}.

\begin{quote}
Armenta bovum in campīs sunt.
\end{quote}

\begin{quote}
\emph{Herds of cattle are on the plains.}
\end{quote}

Second declension, neuter: \emph{armentum, armentī n.} Large animals, cattle and horses; \emph{grex} goes with sheep and goats.

\begin{cousinweb}
\emph{armentum} is usually linked with \emph{arō} 'to plough', since Lewis \& Short give its first sense as cattle for ploughing; EtymDB connects it with \emph{arma} instead, so the origin is not settled.
\end{cousinweb}

\begin{grammarlens}[title={What you've built}]
One more word for telling, explaining and describing at B1.
\end{grammarlens}

\subsection*{Your turn}
\begin{itemize}
\raggedright
  \item \emph{Say it:} \emph{armentum, armentī}
  \item \emph{Say it:} \emph{armentum, armentī}, once more
  \item \emph{Say it:} a sentence of your own with \emph{armentum, armentī}
  \item \emph{Recall:} say the Latin for livestock, then the Latin for a flock, then say \emph{armentum, armentī} again
\end{itemize}

\begin{tcolorbox}[breakable,colback=teal!4,colframe=teal!35!black,title={Before you move on}]
What does armentum, armentī mean? (\textquotedblleft{}A herd (of cattle or horses)\textquotedblright{}.) Say it once more.
\end{tcolorbox}
\section[pāscō, pāscere]{pāscō, pāscere — to pasture, to put out to graze; to feed}

\label{lesson:LA-C477-pasco}



\begin{quote}
Before the new one: say the Latin for a flock, then the Latin for a herd.
\end{quote}

\subsection*{You'll want to know: pāscō, pāscere}
\textbf{pāscō, pāscere} — \textquotedblleft{}to pasture, to put out to graze; to feed\textquotedblright{}.

\begin{quote}
Puer caprās in colle pāscit.
\end{quote}

\begin{quote}
\emph{The boy pastures the goats on the hill.}
\end{quote}

\emph{pāscō, pāscere, pāvī, pāstum}. Passive \emph{pāscor}, 'to graze': \emph{ovēs pāscuntur}. \emph{Pāstor} 'shepherd' comes from it.

\begin{cousinweb}
\emph{pāscō} is inherited from PIE \emph{peh\textsubscript{2}-} 'to protect, to feed' (EtymDB); \emph{pāstor} 'shepherd' and \emph{pābulum} 'fodder' come from the same root (Lewis \& Short). English \emph{food} is a Germanic cousin.
\end{cousinweb}

\begin{grammarlens}[title={What you've built}]
One more word for telling, explaining and describing at B1.
\end{grammarlens}

\subsection*{Your turn}
\begin{itemize}
\raggedright
  \item \emph{Say it:} \emph{pāscō, pāscere}
  \item \emph{Say it:} \emph{pāscō, pāscere}, once more
  \item \emph{Say it:} a sentence of your own with \emph{pāscō, pāscere}
  \item \emph{Recall:} say the Latin for a flock, then the Latin for a herd, then say \emph{pāscō, pāscere} again
\end{itemize}

\begin{tcolorbox}[breakable,colback=teal!4,colframe=teal!35!black,title={Before you move on}]
What does pāscō, pāscere mean? (\textquotedblleft{}To pasture, to put out to graze; to feed\textquotedblright{}.) Say it once more.
\end{tcolorbox}
\section[pāscuum, pāscuī]{pāscuum, pāscuī — a pasture, grazing land}

\label{lesson:LA-C477-pascuum}



\begin{quote}
Before the new one: say the Latin for a herd, then the Latin for to pasture.
\end{quote}

\subsection*{You'll want to know: pāscuum, pāscuī}
\textbf{pāscuum, pāscuī} — \textquotedblleft{}a pasture, grazing land\textquotedblright{}.

\begin{quote}
Vēre pecus in pāscua dūcitur.
\end{quote}

\begin{quote}
\emph{In spring the livestock is taken to the pastures.}
\end{quote}

Second declension, neuter: \emph{pāscuum, pāscuī n.}, mostly plural. The neuter of the adjective \emph{pāscuus} 'for grazing'.

\begin{cousinweb}
\emph{pāscuum} is the neuter of the adjective \emph{pāscuus} 'for grazing', built on \emph{pāscō} 'to pasture' (Lewis \& Short).
\end{cousinweb}

\begin{grammarlens}[title={What you've built}]
One more word for telling, explaining and describing at B1.
\end{grammarlens}

\subsection*{Your turn}
\begin{itemize}
\raggedright
  \item \emph{Say it:} \emph{pāscuum, pāscuī}
  \item \emph{Say it:} \emph{pāscuum, pāscuī}, once more
  \item \emph{Say it:} a sentence of your own with \emph{pāscuum, pāscuī}
  \item \emph{Recall:} say the Latin for a herd, then the Latin for to pasture, then say \emph{pāscuum, pāscuī} again
\end{itemize}

\begin{tcolorbox}[breakable,colback=teal!4,colframe=teal!35!black,title={Before you move on}]
What does pāscuum, pāscuī mean? (\textquotedblleft{}A pasture, grazing land\textquotedblright{}.) Say it once more.
\end{tcolorbox}
